\section{ТЕОРЕТИЧЕСКИЙ АНАЛИЗ МЕТОДОВ ОБНАРУЖЕНИЯ АНОМАЛИЙ И ОБОСНОВАНИЕ ТРЕБОВАНИЙ К СИСТЕМЕ}
\label{ch:theory}

\noindent
Содержание главы соответствует первому этапу постановки задачи:
изучение литературы и существующих методов обнаружения аномалий
в сетевом трафике; сравнительный анализ нейросетевых архитектур
и публично доступных датасетов NIDS на основе балльной оценки с
явно сформулированными критериями и весами; систематизация
подходов к активному тестированию NIDS; формулировка гипотезы
исследования и требований к системе.

\subsection{Корпоративная сеть как объект мониторинга}

Под корпоративной сетью в настоящей работе понимается совокупность
взаимосвязанных программно-аппаратных средств, обеспечивающих обмен
данными между рабочими станциями, серверами, сетевым оборудованием
и внешними сегментами. С точки зрения мониторинга безопасности
корпоративная сеть характеризуется тремя свойствами, существенными
для построения NIDS:

\begin{enumerate}
    \item высокой неоднородностью трафика, приводящей к
    мультимодальному распределению признаков и затрудняющей
    применение единых пороговых правил~\cite{khraisat2019survey,
    ahmad2021nids};
    \item доминированием легитимного трафика, что создает сильный
    дисбаланс классов и требует специализированных методов обучения
    и оценки~\cite{sommer2010outside, chawla2002smote, lin2017focal};
    \item динамичностью характеристик, приводящей к явлению дрейфа
    распределений~\cite{gama2014drift, andresini2021insomnia}.
\end{enumerate}

Объектами мониторинга в работе принимаются: сетевые соединения
(flow) на уровне L3/L4, сессии прикладного уровня и временные
последовательности соединений в скользящих окнах. Это согласуется
с подходом, принятым в большинстве современных систем сетевого
анализа трафика и в открытых датасетах~\cite{moustafa2015unsw,
sharafaldin2018toward}.

В качестве рабочей таксономии угроз использована классификация
датасета UNSW-NB15, выделяющая девять семейств атакующего трафика:
Fuzzers, Analysis, Backdoors, DoS, Exploits, Generic, Reconnaissance,
Shellcode, Worms~\cite{moustafa2015unsw}.

\subsection{Классификация и принципы построения IDS}

В соответствии с устоявшейся в литературе типологией IDS делятся на
три класса~\cite{khraisat2019survey, ahmad2021nids, ferrag2020dlcyber}:
сигнатурные (совпадение с заранее заданным шаблоном), аномальные
(отклонение от модели нормы) и гибридные. Сигнатурный класс
обеспечивает высокую точность по известным угрозам, но не способен
обнаруживать ранее не встречавшиеся атаки. Аномальный класс
расширяет покрытие на неизвестные сценарии за счет повышенной
частоты ложных срабатываний и чувствительности к дрейфу.

С точки зрения исходного представления данных различают
packet-based (анализ полезной нагрузки, DPI) и flow-based анализ.
Стандартами экспорта flow-записей являются Cisco
NetFlow~\cite{rfc3954} и IETF IPFIX~\cite{rfc7011}. В настоящей
работе принят flow-based подход по причине его совместимости с
шифрованным трафиком, ресурсной эффективности и прямой
переносимости на инфраструктуры реальных корпоративных сетей.

\subsection{Методы обнаружения аномалий}

Постановка задачи: для пространства признаков
\(X \subseteq \mathbb{R}^d\) требуется построить функцию скоринга
\(s\colon X \to \mathbb{R}\) и порог \(\tau\) такие, что
\begin{equation}
\hat{y}(x) = \mathbb{I}\bigl[s(x) > \tau\bigr],
\quad x \in X,
\label{eq:detector}
\end{equation}
где \(\hat{y}(x)\) --- предсказанная метка наблюдения,
\(\hat{y}(x)=1\) интерпретируется как <<аномалия/атака>>; \(s(x)\) ---
функция скоринга наблюдения; \(\tau\) --- порог принятия решения;
\(\mathbb{I}[\cdot]\) --- индикаторная функция события.

Статистические и энтропийные методы. Простейший подход ---
многомерное расстояние Махаланобиса:
\begin{equation}
d_M(x) = \sqrt{(x-\mu)^{\top}\Sigma^{-1}(x-\mu)},
\label{eq:mahalanobis}
\end{equation}
где \(\mu\) --- вектор математических ожиданий по обучающей
выборке; \(\Sigma\) --- ковариационная матрица обучающей выборки.

Isolation Forest~\cite{liu2008isolation} определяет скоринг через
ожидаемую глубину изоляции точки:
\begin{equation}
s(x, n) = 2^{-\frac{\mathbb{E}[h(x)]}{c(n)}},\quad
c(n) \approx 2(\ln(n-1) + 0.5772) - \frac{2(n-1)}{n},
\label{eq:iforest}
\end{equation}
где \(n\) --- размер обучающей выборки; \(h(x)\) --- глубина
изоляции наблюдения в случайном дереве; \(c(n)\) --- нормирующий
множитель. Метод имеет линейную сложность и подходит для
unsupervised-режима.

One-Class SVM~\cite{scholkopf2001ocsvm} ищет в спрямляющем
пространстве минимальную гиперсферу, содержащую долю \(1-\nu\)
обучающих наблюдений; параметр \(\nu \in (0,1]\) задает верхнюю
оценку доли выбросов.

\subsection{Сравнительный анализ нейросетевых архитектур}

\subsubsection{Критерии сравнения}

Критерии сравнения нейросетевых архитектур извлечены из
систематических обзоров и привязаны к проверяемым источникам.
Приняты семь критериев C1--C7~\cite{khraisat2019survey,
ahmad2021nids, ferrag2020dlcyber, vinayakumar2019dl}:
C1~--- моделирование локальных паттернов;
C2~--- моделирование временных зависимостей;
C3~--- F1/Accuracy в рецензируемых работах на современных
датасетах NIDS;
C4~--- сложность обучения;
C5~--- inference latency на CPU;
C6~--- зрелость подхода в NIDS (число независимых
воспроизведений);
C7~--- гибкость расширения (embedding, fine-tune, online).

\subsubsection{Краткое описание архитектур}

В работе рассмотрено девять типов нейросетевых архитектур.

MLP реализуется рекурсией
\(h^{(l)} = \phi(W^{(l)} h^{(l-1)} + b^{(l)})\) и служит нижней
границей DL-моделей.

1D-CNN применяет операцию свертки
\((h * w)_t = \sum_{k=0}^{K-1} w_k h_{t-k}\) для извлечения локальных
паттернов в скользящем окне.

LSTM~\cite{hochreiter1997lstm} вводит ячейку состояния \(c_t\) и три
гейта:
\begin{align}
i_t &= \sigma(W_i [h_{t-1}, x_t] + b_i), \nonumber\\
f_t &= \sigma(W_f [h_{t-1}, x_t] + b_f), \nonumber\\
o_t &= \sigma(W_o [h_{t-1}, x_t] + b_o), \nonumber\\
\tilde{c}_t &= \tanh(W_c [h_{t-1}, x_t] + b_c), \nonumber\\
c_t &= f_t \odot c_{t-1} + i_t \odot \tilde{c}_t, \nonumber\\
h_t &= o_t \odot \tanh(c_t),
\label{eq:lstm}
\end{align}
где \(x_t\) --- входной вектор на шаге \(t\);
\(h_{t-1}\) --- скрытое состояние на предыдущем шаге;
\(i_t, f_t, o_t\) --- гейты входа, забывания и выхода;
\(\sigma\) --- сигмоидальная функция активации;
\(\odot\) --- поэлементное произведение.

GRU~\cite{cho2014gru} --- упрощенная RNN с двумя гейтами;
BiLSTM использует прямой и обратный LSTM для двунаправленного
контекста.

Гибрид CNN-LSTM применяет последовательное преобразование
\[
x_t \xrightarrow{\text{Conv1D}} z_t \xrightarrow{\text{LSTM}} h_t
\xrightarrow{\text{FC}} \hat{p}_t,
\]
сочетая локальные паттерны и временные зависимости.

TCN~\cite{bai2018tcn} использует причинные дилатированные свертки с
экспоненциально растущим рецептивным полем.

Transformer-encoder~\cite{vaswani2017attention} основан на
само-внимании:
\begin{equation}
\operatorname{Attention}(Q, K, V) =
\operatorname{softmax}\!\left(\frac{QK^{\top}}{\sqrt{d_k}}\right) V,
\label{eq:attention}
\end{equation}
где \(Q\) --- матрица запросов; \(K\) --- матрица ключей;
\(V\) --- матрица значений; \(d_k\) --- размерность векторов ключей.

Autoencoder и VAE. Автоэнкодер минимизирует ошибку реконструкции
\(\|x - D(E(x))\|^2\); вариационный
автоэнкодер~\cite{kingma2014vae} обучает вероятностное кодирование
максимизацией ELBO:
\begin{equation}
\mathcal{L}_{\text{VAE}} =
\mathbb{E}_{q(z|x)}[\log p(x|z)] - \KL(q(z|x)\|p(z)),
\label{eq:vae}
\end{equation}
где \(x\) --- входное наблюдение; \(z\) --- скрытая (латентная)
переменная; \(q(z|x)\) --- кодировщик; \(p(x|z)\) --- декодировщик;
\(p(z)\) --- априорное распределение латентной переменной;
\(\KL\) --- расхождение Кульбака--Лейблера.

\subsubsection{Балльная оценка и выбор архитектуры}

В таблице~\ref{tab:dl-arch-compare} приведена взвешенная балльная
оценка девяти архитектур по семи критериям. Веса критериев
установлены на основе их значимости для задачи классификации
последовательностей flow-окон в реальном времени с эксплуатационными
ограничениями SOC: C1~--- 0,15; C2~--- 0,20; C3~--- 0,25;
C4~--- 0,05; C5~--- 0,15; C6~--- 0,10; C7~--- 0,10.

\begin{table}[H]
\caption{Балльная оценка нейросетевых архитектур по критериям C1--C7}
\label{tab:dl-arch-compare}
\centering
\small
\begin{tabular}{|L{3.0cm}|c|c|c|c|c|c|c|c|}
\hline
Архитектура & C1 & C2 & C3 & C4 & C5 & C6 & C7 & Балл \\
\hline
MLP                    & 1 & 1 & 2 & 5 & 5 & 2 & 2 & 2,15 \\
\hline
1D-CNN                 & 4 & 2 & 3 & 4 & 5 & 3 & 3 & 3,30 \\
\hline
LSTM                   & 1 & 4 & 3 & 3 & 4 & 4 & 3 & 3,15 \\
\hline
GRU                    & 1 & 4 & 3 & 4 & 4 & 4 & 3 & 3,20 \\
\hline
BiLSTM                 & 2 & 5 & 4 & 2 & 3 & 4 & 3 & 3,55 \\
\hline
TCN                    & 4 & 4 & 3 & 3 & 3 & 3 & 3 & 3,40 \\
\hline
CNN-LSTM      & 5 & 5 & 5 &
                       2 & 4 & 5 &
                       5 & 4,60 \\
\hline
Transformer-encoder    & 3 & 4 & 4 & 2 & 3 & 3 & 4 & 3,45 \\
\hline
AE / VAE               & 2 & 1 & 2 & 3 & 4 & 3 & 3 & 2,40 \\
\hline
\end{tabular}
\end{table}

\noindent
Шкала: 1 --- не соответствует; 5 --- выраженно соответствует.

Лидер по взвешенной балльной оценке --- гибрид CNN-LSTM с баллом
4,60, опережающий ближайших конкурентов BiLSTM (3,55),
Transformer-encoder (3,45) и TCN (3,40). Это решение согласуется с
серией публикаций по NIDS на UNSW-NB15 в 2024--2025 годах, в
которых CNN-LSTM устойчиво обеспечивает наибольшие F1 и Accuracy в
DL-семействе~\cite{vinayakumar2019dl, ahmad2021nids, ferrag2020dlcyber}.

\subsection{Сравнительный анализ датасетов NIDS}

Критерии оценки датасетов NIDS опираются на эталонный фреймворк
Gharib и соавт. (2016)~\cite{gharib2016evalframework} и обзор
Ring и соавт. (2019). Приняты восемь критериев K1--K8:
K1~--- актуальность модели угроз;
K2~--- репрезентативность нормального трафика;
K3~--- качество разметки;
K4~--- сбалансированность распределения классов;
K5~--- пригодность для глубокого обучения;
K6~--- flow-формат и совместимость с NetFlow/IPFIX;
K7~--- воспроизводимость (открытая лицензия, документация);
K8~--- активность использования в современной литературе.

В таблице~\ref{tab:datasets-compare} приведена 5-балльная оценка
десяти публично доступных датасетов NIDS по критериям K1--K8.

\begin{table}[H]
\caption{Балльная оценка публично доступных датасетов NIDS}
\label{tab:datasets-compare}
\centering
\small
\begin{tabular}{|L{3.3cm}|c|c|c|c|c|c|c|c|}
\hline
Датасет & K1 & K2 & K3 & K4 & K5 & K6 & K7 & K8 \\
\hline
DARPA 1998/1999~\cite{lippmann2000darpa} & 1 & 2 & 3 & 3 & 1 & 2 & 4 & 1 \\
\hline
KDDCup99~\cite{tavallaee2009nslkdd}      & 1 & 1 & 2 & 3 & 3 & 2 & 4 & 2 \\
\hline
NSL-KDD~\cite{tavallaee2009nslkdd}       & 1 & 2 & 3 & 4 & 2 & 2 & 5 & 4 \\
\hline
UNSW-NB15~\cite{moustafa2015unsw} & 5 &
5 & 5 & 4 & 5 & 5 &
5 & 5 \\
\hline
CICIDS2017~\cite{sharafaldin2018toward}  & 4 & 4 & 2 & 3 & 4 & 5 & 4 & 5 \\
\hline
CSE-CIC-IDS2018                          & 4 & 4 & 2 & 3 & 4 & 5 & 3 & 4 \\
\hline
TON-IoT~\cite{moustafa2021toniot}        & 4 & 3 & 4 & 3 & 4 & 5 & 4 & 4 \\
\hline
Bot-IoT~\cite{koroniotis2019botiot}      & 3 & 1 & 4 & 1 & 4 & 5 & 4 & 3 \\
\hline
Kyoto 2006+~\cite{song2011kyoto}         & 2 & 2 & 3 & 3 & 3 & 3 & 3 & 2 \\
\hline
LITNET-2020~\cite{damasevicius2020litnet}& 4 & 4 & 3 & 3 & 4 & 4 & 3 & 3 \\
\hline
\end{tabular}
\end{table}

Веса критериев согласованы с приоритетами Gharib и соавт. (2016)
и обзорной работы Ring и соавт. (2019): K1, K2, K3, K5~--- по
0,15; K4, K6, K7, K8~--- по 0,10. Взвешенный балл лидеров:
UNSW-NB15 --- 4,80; TON-IoT --- 3,80; CICIDS2017 --- 3,75.

В работе принят UNSW-NB15 как основной датасет, CICIDS2017 ---
как вспомогательный для cross-evaluation в проверке гипотезы
переносимости методики.

\subsection{Систематизация подходов к активному тестированию NIDS}

В литературе устойчиво выделяются три семейства подходов к
формированию атакующего трафика для тестирования NIDS:

\begin{enumerate}
    \item Шаблонная (rule/template-based) генерация: атаки задаются
    параметрическими шаблонами потоков; применение --- от классических
    работ DARPA/KDD~\cite{lippmann2000darpa,
    gharib2016evalframework} до открытых инструментариев типа ID2T
    и ConCap.
    \item Генеративные модели: GAN, cGAN, WGAN, VAE для синтеза
    реалистичных flow-векторов и последовательностей~\cite{
    ring2019flowgan, lopez2017cvae, goodfellow2014gan,
    mirza2014cgan, arjovsky2017wgan, gulrajani2017wgangp}.
    \item Adversarial-агенты на обучении с подкреплением: агент
    обучается формировать последовательности действий, обходящих
    детектор~\cite{apruzzese2020adversarial}.
\end{enumerate}

В рамках настоящей работы выбран гибридный шаблонный подход с
конечно-автоматным планировщиком и инжектором дрейфа:
\[
\text{Templates} + \text{FSM-планировщик} + \text{Drift-инжектор}.
\]
Аргументы выбора: высокая воспроизводимость и прозрачная
ground-truth; семантическая согласованность шаблонов с таксономией
UNSW-NB15; отсутствие риска коллапса режима, характерного для
GAN-семейства; интерпретируемость для аналитика SOC.

\subsection{Метрики качества}

Для бинарной задачи введены матрица ошибок (TP, TN, FP, FN) и
производные метрики:
\begin{align}
\text{Precision} &= \frac{TP}{TP + FP},
\quad
\text{Recall} = \frac{TP}{TP + FN},
\label{eq:pr}\\
\text{F1} &= \frac{2\,\text{Precision}\,\text{Recall}}
                 {\text{Precision} + \text{Recall}},
\quad
\text{FPR} = \frac{FP}{FP + TN}.\label{eq:f1}
\end{align}

MCC (Matthews correlation coefficient) рекомендуется при сильном
дисбалансе:
\begin{equation}
\text{MCC} =
\frac{TP \cdot TN - FP \cdot FN}
{\sqrt{(TP+FP)(TP+FN)(TN+FP)(TN+FN)}}.
\label{eq:mcc}
\end{equation}

ROC-AUC и PR-AUC интегрируют качество модели по всем порогам;
PR-AUC более чувствителен к редкому классу и принят как основная
интегральная метрика. К эксплуатационным метрикам относятся
time-to-detect (TTD), inference latency, throughput и FP-per-time.

Температурное масштабирование~\cite{guo2017calibration} согласует
выходную вероятность с эмпирической частотой:
\begin{equation}
\hat{p}_{T}(y=1\mid x) = \sigma\!\left(\frac{z(x)}{T}\right),
\label{eq:temp-scaling}
\end{equation}
где \(z(x)\) --- логит, \(T\) --- параметр, подбираемый на
валидационной выборке. Качество калибровки оценивается по ECE
(Expected Calibration Error).

Метрики дрейфа. Population Stability Index, расхождение Кульбака--
Лейблера и Maximum Mean Discrepancy:
\begin{align}
\PSI &= \sum_{b}(p_b - q_b)\ln\frac{p_b}{q_b},
\label{eq:psi}\\
\MMD^{2}(P, Q) &= \mathbb{E}_{P,P}[k(x,x')] - 2\mathbb{E}_{P,Q}[k(x,y)]
                + \mathbb{E}_{Q,Q}[k(y,y')],\label{eq:mmd}
\end{align}
где \(b\) --- индекс интервала разбиения признакового пространства;
\(p_b\) и \(q_b\) --- доли наблюдений в \(b\)-м интервале
эталонной и текущей выборки соответственно;
\(k(\cdot,\cdot)\) --- характеристическое ядро (гауссово ядро с
автобандвидтом по медианной эвристике).

\subsection{Формализация требований и гипотеза}

Концептуальная модель системы состоит из трех подсистем:
предобработки, детекции и активного тестирования
(конечно-автоматный агент). Сформированы 10 функциональных и
10 нефункциональных требований; ключевые количественные пороги
приведены в таблице~\ref{tab:requirements}.

\begin{table}[H]
\caption{Ключевые нефункциональные требования к системе}
\label{tab:requirements}
\centering
\begin{tabularx}{\linewidth}{|c|Z|c|}
\hline
ID & Содержание & Порог \\
\hline
NFR-1a & F1 на UNSW-NB15 (binary) & $\geq 0{,}90$ \\
\hline
NFR-1b & PR-AUC на UNSW-NB15 & $\geq 0{,}95$ \\
\hline
NFR-2 & FPR при Recall $\geq 0{,}85$ & $\leq 0{,}01$ \\
\hline
NFR-3 & Inference latency (CPU, batch=1) & $\leq 50$ мс/окно \\
\hline
NFR-4 & Throughput (CPU, батч) & $\geq 1000$ EPS \\
\hline
NFR-5 & Time-to-detect (DDoS / Scan / Brute) & $\leq 1$ окна \\
\hline
NFR-6 & $\Delta\text{F1}$ при PSI $\leq 0{,}25$ & $\leq 0{,}10$ \\
\hline
NFR-7 & Перенос UNSW $\to$ CICIDS2017 ($\Delta\text{F1}$) & $\leq 0{,}20$ \\
\hline
NFR-8 & Воспроизводимость одной командой & да \\
\hline
NFR-9 & $\sigma_{\text{F1}}$ при $n \geq 3$ seed & $\leq 0{,}01$ \\
\hline
\end{tabularx}
\end{table}

Формально сформулирована гипотеза исследования: применение
нейронных сетей, обрабатывающих последовательности агрегированных
признаков сетевых потоков в скользящих окнах, в сочетании с
программным агентом на основе конечного автомата с библиотекой
параметризованных шаблонов и контролируемым инжектированием дрейфа
распределений, обеспечивает построение методики обнаружения,
превосходящей классические сигнатурные и статистические подходы по
комплексному критерию качества (точность, полнота, FPR при заданной
полноте, задержка детекции, устойчивость к дрейфу). Проверка
проводится на этапе~\ref{ch:experiments}.

\noindent Выводы по главе 1

\begin{enumerate}
    \item Корпоративная сеть как объект мониторинга характеризуется
    неоднородностью трафика, дисбалансом классов и динамичностью
    распределений; обоснован выбор flow-уровня анализа,
    совместимого с NetFlow/IPFIX.
    \item Сигнатурные, аномальные и гибридные подходы NIDS
    дополняют, а не заменяют друг друга; в работе принят аномальный
    нейросетевой подход.
    \item Рассмотрены девять типов нейросетевых архитектур; по
    результатам взвешенной балльной оценки лидер --- гибрид
    CNN-LSTM с баллом 4,60, принимаемый в качестве основной
    proposed-архитектуры.
    \item Из десяти публично доступных датасетов NIDS лидером по
    взвешенной балльной оценке является UNSW-NB15 (балл 4,80);
    как вспомогательный принят CICIDS2017.
    \item Систематизированы три семейства подходов к моделированию
    атак; обоснован выбор гибридного шаблонного подхода с
    конечно-автоматным планировщиком без cGAN-составляющей.
    \item Сформированы количественные требования (10 функциональных,
    10 нефункциональных) и гипотеза исследования.
\end{enumerate}
